\chapter{نقشهٔ راه اجرایی}
\label{ch:roadmap}
%=====================================================================

\section{منطق فازبندی}

نقشهٔ راهی که در این فصل ارائه می‌شود، بر یک اصل ساده بنا شده است: در پروژه‌ای که بخش قابل‌توجهی از تجهیزات حیاتی آن (بخش \ref{sec:labequipment}) به‌دلیل تحریم فناوری با عدم‌قطعیت زمان تحویل مواجه است، بزرگ‌ترین خطر برای زمان‌بندی، پیچیدگی مهندسی نیست بلکه \textbf{زنجیرهٔ تأمین} است. به همین دلیل، برخلاف الگوی متداول «ابتدا طراحی کامل شود، سپس خرید انجام گیرد، سپس ساخته شود»، در این طرح فاز تدارکات اقلام حساس از همان ماه‌های نخست و به‌موازات فازهای تحقیق و طراحی آغاز می‌شود. اگر این همپوشانی رعایت نشود، تیم پروژه می‌تواند در پایان ماه ششم طراحی کاملی در دست داشته باشد اما همچنان چند ماه دیگر منتظر رسیدن یک آشکارساز تک‌فوتون بماند -- سناریویی که کاملاً قابل‌پیش‌بینی و قابل‌اجتناب است.

فازبندی به نه بخش اصلی تقسیم شده که سه مورد از آن‌ها (فاز \lr{0}، فاز \lr{2} و فاز \lr{6}) به‌صورت موازی با بقیه پیش می‌روند نه به‌صورت متوالی؛ توضیح هر فاز، همراه با نتیجهٔ ملموسی که در پایان آن باید روی میز باشد، در ادامه آمده است.

\subsection*{فاز \texorpdfstring{\lr{0}}{۰} -- تحقیق و توسعه (ماه \texorpdfstring{\lr{1}}{۱} تا \texorpdfstring{\lr{3}}{۳})}
این فاز، پایهٔ فکری کل پروژه است و شامل مرور نظام‌مند ادبیات فنی روز (از جمله منابع فصل‌های \ref{chap:hybrid} و بخش \ref{sec:pqcfoundations})، انتخاب نهایی پروتکل \lr{QKD} (بین \lr{decoy-state BB84}، \lr{MDI-QKD} و \lr{TF-QKD}) متناسب با فاصلهٔ واقعی گره‌های تهران، انتخاب مجموعهٔ الگوریتم‌های \lr{PQC} استاندارد \lr{NIST} (\lr{ML-KEM}، \lr{ML-DSA}، \lr{SLH-DSA} به‌عنوان پشتیبان)، و شبیه‌سازی نرم‌افزاری کامل کانال و پروتکل پیش از هرگونه ساخت سخت‌افزاری است. خروجی این فاز یک سند فنی مرجع است که مبنای تصمیم‌گیری تمام فازهای بعدی قرار می‌گیرد.

\subsection*{فاز \texorpdfstring{\lr{1}}{۱} -- طراحی شبکه (ماه \texorpdfstring{\lr{2}}{۲} تا \texorpdfstring{\lr{4}}{۴})}
موازی با پایان‌یافتن فاز \lr{0}، تیم شبکه باید نقشه‌برداری فیبر تاریک موجود مخابرات را انجام دهد، تلفات واقعی مسیرهای پیشنهادی را با اندازه‌گیری میدانی (نه صرفاً برآورد نظری) تعیین کند، گره‌های نهایی را قطعی و توپولوژی حلقه‌ای شش‌گانه را نهایی نماید، و سند الزامات امنیت فیزیکی سایت‌ها را تدوین کند.

\subsection*{فاز \texorpdfstring{\lr{2}}{۲} -- تدارکات و تأمین تجهیزات (ماه \texorpdfstring{\lr{2}}{۲} تا \texorpdfstring{\lr{8}}{۸}، موازی)}
مهم‌ترین فاز از منظر مدیریت ریسک. سفارش اقلام کم‌ریسک (سرور، شبکه، برد \lr{FPGA} میان‌رده) می‌تواند در همان هفته‌های نخست و با زمان تحویل کوتاه انجام شود؛ اما سفارش اقلام پرریسک تحریمی -- آشکارساز تک‌فوتون، مدولاتور نوری، \lr{QRNG} سخت‌افزاری (بخش \ref{sec:labequipment}) -- باید همزمان و بدون تأخیر آغاز شود، چون این دقیقاً همان اقلامی هستند که زمان تحویل چندماهه و غیرقطعی دارند. توصیهٔ صریح این نقشهٔ راه آن است که مسیر تأمین این اقلام حساس، به‌موازات مسیر توسعهٔ داخلی جایگزین (همکاری با مراکز پژوهشی اپتوالکترونیک داخلی) دنبال شود تا در صورت تأخیر یا شکست یک مسیر، پروژه به‌طور کامل متوقف نشود.

\subsection*{فاز \texorpdfstring{\lr{3}}{۳} -- نمونهٔ اولیهٔ آزمایشگاهی \texorpdfstring{\lr{QKD}}{QKD} (ماه \texorpdfstring{\lr{4}}{۴} تا \texorpdfstring{\lr{8}}{۸})}
با رسیدن نخستین محمولهٔ تجهیزات، سامانهٔ فرستنده/گیرنده روی میز آزمایشگاه (نه در فیبر میدانی) مونتاژ و کالیبره می‌شود: هم‌ترازی نوری، کالیبراسیون آشکارساز، و اندازه‌گیری نرخ خطای بیت کوانتومی (\lr{QBER}) در یک حلقهٔ فیبر کوتاه و کاملاً کنترل‌شده. این فاز عمداً پیش از هرگونه استقرار میدانی قرار گرفته تا خطاهای اساسی طراحی، پیش از صرف هزینهٔ حفاری و نصب فیبر شهری کشف شوند.

\subsection*{فاز \texorpdfstring{\lr{4}}{۴} -- پیوند آزمایشی میدانی دوگانه (ماه \texorpdfstring{\lr{7}}{۷} تا \texorpdfstring{\lr{10}}{۱۰})}
انتقال سامانهٔ اثبات‌شده از آزمایشگاه به نخستین پیوند واقعی میان دو گرهٔ اولویت‌دار (مثلاً بانک مرکزی و مرکز دادهٔ شتاب/شاپرک)، به‌عنوان اثبات مفهوم روی زیرساخت واقعی مخابراتی تهران.

\subsection*{فاز \texorpdfstring{\lr{5}}{۵} -- تکمیل حلقهٔ کلان‌شهری (ماه \texorpdfstring{\lr{9}}{۹} تا \texorpdfstring{\lr{13}}{۱۳})}
افزودن تدریجی باقی گره‌ها و تکمیل توپولوژی حلقه‌ای شش‌گانه، به‌همراه آزمون تاب‌آوری شبکه در برابر قطعی یک پیوند منفرد.

\subsection*{فاز \texorpdfstring{\lr{6}}{۶} -- لایهٔ نرم‌افزاری \texorpdfstring{\lr{PQC}}{PQC} (ماه \texorpdfstring{\lr{2}}{۲} تا \texorpdfstring{\lr{9}}{۹}، موازی)}
از بزرگ‌ترین مزیت‌های این طرح آن است که یکپارچه‌سازی \lr{ML-KEM/ML-DSA} در نرم‌افزار بانکی، کاملاً مستقل از سخت‌افزار فیزیکی \lr{QKD} پیش می‌رود و به هیچ‌یک از تأخیرهای احتمالی فازهای \lr{2} تا \lr{5} وابسته نیست. این استقلال به‌معنای آن است که حتی اگر تأمین سخت‌افزار کوانتومی با تأخیر مواجه شود، ارتقای امنیت رمزنگاری کلاسیک بانک همچنان طبق برنامه پیش می‌رود -- و از این جهت، فاز \lr{6} «شبکهٔ اطمینان» کل پروژه در برابر ریسک تحریم است.

\subsection*{فاز \texorpdfstring{\lr{7}}{۷} -- یکپارچه‌سازی کامل بانکی (ماه \texorpdfstring{\lr{10}}{۱۰} تا \texorpdfstring{\lr{14}}{۱۴})}
اتصال \lr{KMS} به سامانه‌های پیام‌رسانی مالی موجود و آزمون تراکنش‌های واقعی در محیط کنترل‌شده (\lr{sandbox})، پیش از هرگونه اتصال به محیط عملیاتی زندهٔ بانک.

\subsection*{فاز \texorpdfstring{\lr{8}}{۸} -- ممیزی امنیتی و بهره‌برداری پایدار (ماه \texorpdfstring{\lr{13}}{۱۳} تا \texorpdfstring{\lr{16}}{۱۶} و ادامه)}
انتقال به بهره‌برداری کامل، عبور از یک دور آزمون نفوذ مستقل، و آغاز برنامه‌ریزی برای توسعهٔ آتی به مقیاس بین‌شهری (بخش \ref{sec:pqcfoundations}، فصل \ref{chap:cost}).

\begin{figure}[htbp]
\centering
\begin{ganttchart}[
  hgrid,vgrid,
  bar/.append style={fill=blue!45},
  bar height=0.55,
  title height=1,
  x unit=0.62cm,
  y unit chart=0.62cm,
]{1}{16}
\gantttitle{بازهٔ زمانی (ماه)}{16} \\
\gantttitlelist{1,...,16}{1} \\
\ganttbar{فاز \lr{0}: تحقیق و توسعه}{1}{3} \\
\ganttbar{فاز \lr{1}: طراحی شبکه}{2}{4} \\
\ganttbar[bar/.append style={fill=red!45}]{فاز \lr{2}: تدارکات (پرریسک تحریمی)}{2}{8} \\
\ganttbar{فاز \lr{3}: نمونهٔ اولیهٔ آزمایشگاهی}{4}{8} \\
\ganttbar{فاز \lr{4}: پیوند آزمایشی میدانی}{7}{10} \\
\ganttbar{فاز \lr{5}: تکمیل حلقهٔ کلان‌شهری}{9}{13} \\
\ganttbar[bar/.append style={fill=green!45}]{فاز \lr{6}: لایهٔ نرم‌افزاری \lr{PQC}}{2}{9} \\
\ganttbar{فاز \lr{7}: یکپارچه‌سازی بانکی}{10}{14} \\
\ganttbar{فاز \lr{8}: ممیزی و بهره‌برداری}{13}{16}
\end{ganttchart}
\caption{نقشهٔ راه تفصیلی اجرای فاز پایلوت؛ نوار قرمز (فاز \lr{2}) گلوگاه اصلی ریسک تحریمی و نوار سبز (فاز \lr{6}) مسیر مستقل و کم‌ریسک \lr{PQC} نرم‌افزاری را نشان می‌دهد}
\label{fig:gantt}
\end{figure}

\section{دو سناریوی زمانی: دست‌بالا در برابر واقع‌بینانه}

از آنجا که مدت زمان فاز \lr{2} (تدارکات) بیش از هر عامل دیگری به شرایط بیرونی و غیرقابل‌کنترل تیم پروژه وابسته است، ارائهٔ یک عدد قطعی برای «طول کل پروژه» گمراه‌کننده خواهد بود. جدول \ref{tab:scenario} دو سناریوی صریح را در کنار هم می‌گذارد تا تصمیم‌گیرندگان بتوانند بودجه و انتظارات را بر مبنای واقعیت تنظیم کنند، نه بر مبنای خوش‌بینانه‌ترین حالت ممکن.

\begin{table}[htbp]
\centering
\caption{مقایسهٔ سناریوی دست‌بالا (خوش‌بینانه) و واقع‌بینانه برای افق زمانی طرح پایلوت}
\label{tab:scenario}
\small
\begin{tabular}{p{2.8cm}p{2.2cm}p{8.5cm}}
\toprule
\textbf{سناریو} & \textbf{افق کل} & \textbf{مفروضات کلیدی} \\
\midrule
دست‌بالا (خوش‌بینانه) & $12$ ماه & تأمین اقلام حساس تحریمی بدون تأخیر و در همان بازهٔ برنامه‌ریزی‌شده (ماه $2$ تا $6$) انجام شود؛ هیچ‌یک از گره‌های فیبر نیاز به حفاری تکمیلی نداشته باشند؛ تخصیص بودجه از همان ماه اول قطعی و بدون وقفهٔ اداری باشد. \\
واقع‌بینانه & $16$ تا $18$ ماه & فاز تدارکات اقلام حساس ($2$ تا $3$ قلم اصلی) با میانگین $3$ تا $5$ ماه تأخیر نسبت به برنامه مواجه شود؛ حداقل یک دور بازطراحی پس از آزمون آزمایشگاهی (فاز $3$) لازم باشد؛ فرآیندهای اداری تخصیص بودجه و مناقصهٔ دولتی $1$ تا $2$ ماه به زمان اولیه اضافه کنند. \\
\bottomrule
\end{tabular}
\end{table}

توصیهٔ این نقشهٔ راه آن است که در هرگونه سند رسمی ارائه‌شده به نهادهای تصمیم‌گیرنده، عدد \textbf{واقع‌بینانه} (نه دست‌بالا) به‌عنوان افق رسمی پروژه اعلام شود؛ اعلام عدد خوش‌بینانه به‌عنوان تعهد رسمی، رایج‌ترین علت شکست اعتبار پروژه‌های زیرساختی در برخورد اول با تأخیر اجتناب‌ناپذیر یک قلم تحریمی است.

\section{شاخص‌های کلیدی موفقیت \lr{(KPI)}}

\begin{itemize}
\item برقراری پایدار پیوند \lr{QKD} با نرخ خطای بیت کوانتومی کمتر از آستانهٔ امنیتی (رابطهٔ \ref{eq:qber}) به‌مدت حداقل $90$ روز متوالی؛
\item یکپارچه‌سازی موفق \lr{PQC (ML-KEM/ML-DSA)} در حداقل یک سامانهٔ پیام‌رسانی مالی داخلی؛
\item عبور موفق از یک دور آزمون نفوذ مستقل بدون یافتن آسیب‌پذیری بحرانی؛
\item تربیت حداقل $10$ متخصص داخلی در حوزهٔ طراحی، بهره‌برداری و ممیزی امنیت \lr{QKD/PQC}؛
\item دستیابی به حداقل یک منبع تأمین داخلی (حتی در سطح نمونهٔ اولیه) برای دست‌کم یکی از اقلام پرریسک تحریمی جدول \ref{tab:labgear}، به‌عنوان گام نخست کاهش وابستگی بلندمدت.
\end{itemize}

%=====================================================================
